\documentclass[11pt,a4paper]{article}

% Paquetes
\usepackage[utf8]{inputenc}
\usepackage[T1]{fontenc}
\usepackage[spanish]{babel}
\usepackage[margin=2.5cm]{geometry}
\usepackage{graphicx}
\usepackage{xcolor}
\usepackage{colortbl}
\usepackage{booktabs}
\usepackage{longtable}
\usepackage{array}
\usepackage{enumitem}
\usepackage{titlesec}
\usepackage{fancyhdr}
\usepackage{tcolorbox}
\usepackage{hyperref}
\usepackage{listings}
\usepackage{multicol}
\usepackage{tikz}
\usetikzlibrary{shapes, arrows.meta, positioning}

% Colores
\definecolor{primary}{HTML}{2563EB}
\definecolor{critical}{HTML}{DC2626}
\definecolor{success}{HTML}{10B981}
\definecolor{warning}{HTML}{F59E0B}
\definecolor{codebg}{HTML}{F1F5F9}
\definecolor{codeframe}{HTML}{CBD5E1}
\definecolor{bashbg}{HTML}{1E293B}
\definecolor{awsorange}{HTML}{FF9900}
\definecolor{ocired}{HTML}{C74634}

% Hyperlinks
\hypersetup{
    colorlinks=true,
    linkcolor=primary,
    urlcolor=primary,
    citecolor=primary,
    pdftitle={Guia de Implementacion - Sentinel AcademIA},
    pdfauthor={Equipo Sentinel AcademIA}
}

% Listings - definicion de lenguajes
\lstdefinelanguage{Python}{
  keywords={async, await, def, return, if, else, elif, for, while, try, except, finally, raise, with, as, from, import, class, lambda, pass, break, continue, in, is, not, and, or, None, True, False, yield, global, nonlocal},
  keywordstyle=\color{primary}\bfseries,
  ndkeywords={str, int, float, bool, list, dict, tuple, set, Optional, Any, Dict, List, Union, Tuple, Set},
  ndkeywordstyle=\color{warning}\bfseries,
  sensitive=true,
  comment=[l]{\#},
  morecomment=[s]{"""}{"""},
  commentstyle=\color{gray}\itshape,
  stringstyle=\color{success},
  morestring=[b]',
  morestring=[b]"
}

\lstdefinelanguage{Bash}{
  keywords={sudo, apt-get, dnf, yum, brew, npm, npx, pip, sam, aws, oci, cd, ls, mkdir, cat, echo, export, source, curl, wget, git, python, node, terraform, kubectl, docker, vim, nano, less, more, grep, awk, sed},
  keywordstyle=\color{success}\bfseries,
  ndkeywords={--version, --help, --region, --profile, --bucket, --stack-name, --guided, --capabilities, --s3-bucket, --no-fail-on-empty-changeset, --no-confirm-changeset, --parameter-overrides},
  ndkeywordstyle=\color{warning}\bfseries,
  sensitive=true,
  comment=[l]{\#},
  commentstyle=\color{gray}\itshape,
  stringstyle=\color{primary},
  morestring=[b]",
  morestring=[b]'
}

\lstdefinelanguage{YAML}{
  keywords={true, false, null, yes, no},
  keywordstyle=\color{primary}\bfseries,
  sensitive=true,
  comment=[l]{\#},
  commentstyle=\color{gray}\itshape,
  stringstyle=\color{success},
  morestring=[b]",
  morestring=[b]'
}

\lstdefinelanguage{JSON}{
  keywords={true, false, null},
  keywordstyle=\color{primary}\bfseries,
  commentstyle=\color{gray}\itshape,
  stringstyle=\color{success},
  morestring=[b]",
  showstringspaces=false
}

\lstset{
    basicstyle=\ttfamily\small,
    backgroundcolor=\color{codebg},
    frame=single,
    framerule=0.5pt,
    rulecolor=\color{codeframe},
    breaklines=true,
    showstringspaces=false,
    numbers=left,
    numberstyle=\tiny\color{gray},
    numbersep=8pt,
}

% Formato titulos
\titleformat{\section}{\Large\bfseries\color{primary}}{\thesection}{1em}{}
\titleformat{\subsection}{\large\bfseries\color{primary!80!black}}{\thesubsection}{1em}{}
\titleformat{\subsubsection}{\normalsize\bfseries\color{primary!70!black}}{\thesubsubsection}{1em}{}

% Header/footer
\pagestyle{fancy}
\fancyhf{}
\rhead{\small Sentinel AcademIA \textbar{} Guia de Implementacion}
\lhead{\small \leftmark}
\rfoot{\thepage}
\renewcommand{\headrulewidth}{0.4pt}

% Title
\title{
    \vspace{-1.5cm}
    {\Huge\bfseries\color{primary} Gu\'ia de Implementaci\'on}\\[0.4cm]
    {\Large Sentinel AcademIA -- Paso a Paso}\\[0.6cm]
    {\normalsize Backend Python + Frontend Vue 3 + AWS + Oracle OCI}
}
\author{Equipo de Desarrollo}
\date{\today}

\begin{document}

\maketitle
\thispagestyle{fancy}

\begin{tcolorbox}[colback=codebg, colframe=primary, title=\textbf{Sobre esta gu\'ia}]
Esta es una gu\'ia \textbf{paso a paso} para implementar Sentinel AcademIA. Est\'a escrita asumiendo que \textbf{NO tienes asistencia} de una IA, por lo que cada paso incluye:

\begin{itemize}[leftmargin=*]
    \item \textbf{Qu\'e} estamos haciendo y por qu\'e
    \item \textbf{C\'omo} hacerlo (comandos exactos)
    \item \textbf{Output esperado} (qu\'e deber\'ias ver)
    \item \textbf{Problemas comunes} y c\'omo resolverlos
\end{itemize}

\textbf{Tiempo estimado total}: 14-18 horas de desarrollo concentrado.
\end{tcolorbox}

\tableofcontents
\newpage

% ============================================
\section{Pre-requisitos}
% ============================================

\subsection{Cuentas necesarias}

\begin{longtable}{|p{4cm}|p{4cm}|p{6cm}|}
\hline
\rowcolor{primary!20}
\textbf{Cuenta} & \textbf{Costo} & \textbf{Para qu\'e} \\
\hline
\endhead
AWS Academy & \$50 USD cr\'editos (gratis) & Lambda, SQS, SNS, DynamoDB, S3, SES, CloudWatch \\
\hline
Oracle Cloud Academy & Cr\'editos disponibles & OCI Generative AI (LLM), Document AI, Vault \\
\hline
Vercel & Free tier & Deploy del frontend Vue 3 \\
\hline
GitHub & Free para public & Repositorio + CI/CD con GitHub Actions \\
\hline
OpenAI/Anthropic (opcional) & - & Backup para LLMs en caso de emergencia \\
\hline
\end{longtable}

\subsection{Herramientas a instalar localmente}

\begin{lstlisting}[language=Bash, caption={Instalacion de herramientas en Linux (Fedora/RHEL)}]
# Actualizar sistema
sudo dnf update -y

# Python 3.12
sudo dnf install -y python3.12 python3.12-pip python3.12-devel

# Node.js 20
curl -fsSL https://rpm.nodesource.com/setup_20.x | sudo bash -
sudo dnf install -y nodejs

# AWS CLI v2
curl "https://awscli.amazonaws.com/awscli-exe-linux-x86_64.zip" -o "awscliv2.zip"
unzip awscliv2.zip
sudo ./aws/install

# SAM CLI
pip3.12 install --user aws-sam-cli

# OCI CLI
bash -c "$(curl -L https://raw.githubusercontent.com/oracle/oci-cli/master/scripts/install/install.sh)"

# Tectonic (compilador LaTeX)
curl -fsSL "https://github.com/tectonic-typesetting/tectonic/releases/download/tectonic%400.16.9/tectonic-0.16.9-x86_64-unknown-linux-musl.tar.xz" -o /tmp/tectonic.tar.xz
tar -xJf /tmp/tectonic.tar.xz -C /tmp
sudo mv /tmp/tectonic /usr/local/bin/

# Git, jq, make
sudo dnf install -y git jq make

# Verificar
python3.12 --version
node --version
aws --version
sam --version
oci --version
tectonic --version
git --version
\end{lstlisting}

\begin{lstlisting}[language=Bash, caption={Instalacion en macOS con Homebrew}]
brew install python@3.12 node awscli aws-sam-cli oci-cli tectonic git jq
brew link --overwrite python@3.12
\end{lstlisting}

\begin{lstlisting}[language=Bash, caption={Instalacion en Windows}]
# Descargar installers desde:
# - Python: https://www.python.org/downloads/
# - Node.js: https://nodejs.org/
# - AWS CLI: https://awscli.amazonaws.com/AWSCLIV2.msi
# - SAM CLI: https://github.com/aws/aws-sam-cli/releases
# - OCI CLI: https://docs.oracle.com/en-us/iaas/Content/API/SDKDocs/cliinstall.htm
\end{lstlisting}

% ============================================
\section{FASE 0: Setup Inicial}
% ============================================

\subsection{0.1 Configurar AWS CLI}

\subsubsection{Obtener credenciales de AWS Academy}

\begin{enumerate}[leftmargin=*]
    \item Ve a tu laboratorio en AWS Academy
    \item Click en \textbf{``AWS Details''} (esquina superior derecha)
    \item Click en \textbf{``Show''} junto a AWS CLI credentials
    \item Copia los 3 valores:
    \begin{itemize}
        \item \texttt{AWS\_ACCESS\_KEY\_ID}
        \item \texttt{AWS\_SECRET\_ACCESS\_KEY}
        \item \texttt{AWS\_SESSION\_TOKEN} (este caduca en 3-4 horas)
    \end{itemize}
\end{enumerate}

\subsubsection{Configurar el archivo de credenciales}

\textbf{Nunca} pongas credenciales en c\'odigo o archivos del proyecto. Usa el archivo de credenciales de AWS CLI:

\begin{lstlisting}[language=Bash, caption={Configurar AWS CLI}]
mkdir -p ~/.aws

# Crear archivo de credenciales
cat > ~/.aws/credentials << 'EOF'
[default]
aws_access_key_id = TU_ACCESS_KEY_AQUI
aws_secret_access_key = TU_SECRET_KEY_AQUI
aws_session_token = TU_SESSION_TOKEN_AQUI
EOF

# Crear archivo de configuracion
cat > ~/.aws/config << 'EOF'
[default]
region = us-east-1
output = json
EOF

# Permisos correctos (importante)
chmod 600 ~/.aws/credentials ~/.aws/config
\end{lstlisting}

\textbf{IMPORTANTE}: Las credenciales de AWS Academy \textbf{caducan cada 3-4 horas}. Cuando caduquen, repite el proceso.

\subsubsection{Verificar acceso}

\begin{lstlisting}[language=Bash, caption={Verificar que las credenciales funcionan}]
# Identidad
aws sts get-caller-identity
# Output esperado:
# {
#     "UserId": "AROA...",
#     "Account": "123456789012",
#     "Arn": "arn:aws:sts::123456789012:assumed-role/LabRole/..."
# }

# Verificar region
aws configure get region
# Output: us-east-1

# Verificar que LabRole existe
aws iam get-role --role-name LabRole

# Crear bucket para SAM artifacts
BUCKET_NAME="sentinel-sam-artifacts-$(date +%s)"
aws s3 mb s3://$BUCKET_NAME --region us-east-1
echo "Bucket creado: $BUCKET_NAME"
# GUARDAR este nombre, lo necesitaras despues
echo $BUCKET_NAME > ~/.sentinel-sam-bucket
\end{lstlisting}

\subsection{0.2 Configurar Oracle OCI CLI}

\subsubsection{Crear API key en OCI}

\begin{enumerate}[leftmargin=*]
    \item Ve a OCI Console: \url{https://cloud.oracle.com/}
    \item Click en tu perfil (esquina superior derecha) $\rightarrow$ \textbf{User Settings}
    \item Seccion \textbf{API Keys} $\rightarrow$ Click \textbf{``Add API Key''}
    \item Click \textbf{``Download Private Key''} y \textbf{``Download Public Key''}
    \item Click \textbf{``Add''}
    \item \textbf{Copia el Fingerprint} que aparece (lo necesitaras)
\end{enumerate}

\subsubsection{Configurar OCI CLI}

\begin{lstlisting}[language=Bash, caption={Instalar y configurar OCI CLI}]
# Instalar
bash -c "$(curl -L https://raw.githubusercontent.com/oracle/oci-cli/master/scripts/install/install.sh)"

# Agregar al PATH (agregar a tu ~/.bashrc o ~/.zshrc)
echo 'export PATH="$HOME/lib/oracle-cli/bin:$PATH"' >> ~/.bashrc
source ~/.bashrc

# Verificar instalacion
oci --version
\end{lstlisting}

\begin{lstlisting}[language=Bash, caption={Configurar OCI con el config file}]
# Crear directorio de OCI
mkdir -p ~/.oci

# Mover las keys descargadas
mv ~/Downloads/api_key.pem ~/.oci/
mv ~/Downloads/*_fingerprint ~/.oci/  # el fingerprint file

# Configurar permisos
chmod 600 ~/.oci/api_key.pem

# Crear config file
cat > ~/.oci/config << EOF
[DEFAULT]
user=ocid1.user.oc1..TU_USER_OCID
tenancy=ocid1.tenancy.oc1..TU_TENANCY_OCID
region=us-chicago-1
key_file=$HOME/.oci/api_key.pem
fingerprint=TU_FINGERPRINT_AQUI
EOF

chmod 600 ~/.oci/config
\end{lstlisting}

\begin{lstlisting}[language=Bash, caption={Verificar acceso a OCI}]
# Test basico
oci iam region-subscription list
# Debe listar las regiones disponibles

# Obtener tu Tenancy OCID (si no lo tienes)
oci iam tenancy get --tenancy-id ocid1.tenancy.oc1..TU_TENANCY_OCID
\end{lstlisting}

\subsubsection{Crear Compartment y obtener OCIDs}

\begin{lstlisting}[language=Bash, caption={Setup de OCI}]
# Obtener Tenancy OCID
TENANCY_OCID=$(aws configure get tenancy_id 2>/dev/null || oci iam compartment list --all --query 'data[0]."compartment-id"' --raw-output 2>/dev/null)
# Si no funciona, sacalo del OCI Console

# Crear compartment para el proyecto
oci iam compartment create \
  --compartment-id $TENANCY_OCID \
  --name "SentinelAcademia" \
  --description "Compartment para hackaton Sentinel AcademIA"

# GUARDAR el OCID del compartment que se creo
COMPARTMENT_OCID="ocid1.compartment.oc1..AQUI_TU_COMPARTMENT_OCID"
echo "export OCI_COMPARTMENT_ID=$COMPARTMENT_OCID" >> ~/.bashrc
echo $COMPARTMENT_OCID > ~/.sentinel-compartment
\end{lstlisting}

\subsubsection{Probar OCI Generative AI}

\begin{lstlisting}[language=Bash, caption={Test de OCI GenAI}]
# Verificar modelos disponibles
oci generative-ai-inference model list \
  --compartment-id $COMPARTMENT_OCID \
  --query 'data[*].display-name' \
  --output table

# Probar Cohere Command R
oci generative-ai-inference generate-text \
  --compartment-id $COMPARTMENT_OCID \
  --serving-mode '{"modelId":"cohere.command-r","servingType":"ON_DEMAND"}' \
  --inference-request '{"prompt":"Di hola en espanol","maxTokens":50}'

# Si funciona, veras un JSON con la respuesta
\end{lstlisting}

\subsection{0.3 Clonar el repositorio}

\begin{lstlisting}[language=Bash, caption={Setup del proyecto}]
# Clonar (o descargar como ZIP si es privado)
git clone https://github.com/tu-usuario/sentinel-academia.git
cd sentinel-academia

# Instalar dependencias del backend
cd src
python3.12 -m venv .venv
source .venv/bin/activate
pip install -r requirements.txt
cd ..

# Instalar dependencias del frontend
cd web
npm install
cd ..

# Verificar estructura
ls -la
\end{lstlisting}

% ============================================
\section{FASE 1: Backend Core (Python Lambdas)}
% ============================================

\subsection{1.1 Setup del proyecto Python}

\subsubsection{Estructura de carpetas}

\begin{lstlisting}[language=Bash, caption={Crear estructura del backend}]
cd src
mkdir -p handlers/_shared
mkdir -p services
mkdir -p types
mkdir -p tests/{fixtures,handlers,services}
touch handlers/__init__.py
touch handlers/_shared/__init__.py
touch services/__init__.py
touch types/__init__.py
touch tests/__init__.py
\end{lstlisting}

\subsubsection{Archivo de dependencias (requirements.txt)}

\begin{lstlisting}[caption={src/requirements.txt}]
boto3>=1.34
boto3-stubs[dynamodb,lambda,s3,sqs,sns,events]>=1.34
aws-lambda-powertools[tracer,metrics,validation]>=2.30
pydantic>=2.6
pydantic-settings>=2.2
opossum>=1.4
p-retry>=1.0
rate-limiter-flexible>=2.4
oci>=2.119

# Dev dependencies
pytest>=8.0
pytest-mock>=3.12
pytest-cov>=4.1
moto[dynamodb,s3,sqs,sns,events]>=5.0
freezegun>=1.4
responses>=0.25
ruff>=0.3
mypy>=1.8
black>=24.0
boto3-stubs[essential]>=1.34
\end{lstlisting}

\subsubsection{pyproject.toml}

\begin{lstlisting}[caption={src/pyproject.toml}]
[project]
name = "sentinel-academia-backend"
version = "0.1.0"
requires-python = ">=3.12"
dependencies = [
    "boto3>=1.34",
    "boto3-stubs[dynamodb,lambda,s3,sqs,sns]>=1.34",
    "aws-lambda-powertools[tracer,metrics,validation]>=2.30",
    "pydantic>=2.6",
    "pydantic-settings>=2.2",
    "opossum>=1.4",
    "p-retry>=1.0",
    "rate-limiter-flexible>=2.4",
    "oci>=2.119",
]

[project.optional-dependencies]
dev = [
    "pytest>=8.0",
    "pytest-mock>=3.12",
    "pytest-cov>=4.1",
    "moto[dynamodb,s3,sqs,sns,events]>=5.0",
    "freezegun>=1.4",
    "responses>=0.25",
    "ruff>=0.3",
    "mypy>=1.8",
    "black>=24.0",
]

[tool.ruff]
line-length = 100
target-version = "py312"

[tool.ruff.lint]
select = ["E", "F", "W", "I", "N", "UP", "B", "C4", "DTZ", "RET", "SIM"]

[tool.mypy]
python_version = "3.12"
strict = true
warn_unused_ignores = true
warn_return_any = true

[tool.pytest.ini_options]
testpaths = ["tests"]
addopts = "-v --cov=. --cov-report=term-missing"
\end{lstlisting}

\subsubsection{Variables de entorno (.env)}

\begin{lstlisting}[caption={src/.env}]
# AWS
AWS_REGION=us-east-1
DYNAMODB_TABLE=quejas
S3_BUCKET=sentinel-quejas-adjuntos-DONDE_VA_TU_ACCOUNT_ID
PROCESS_QUEUE_URL=https://sqs.us-east-1.amazonaws.com/DONDE_VA_TU_ACCOUNT_ID/complaints-queue
CRITICAL_ALERTS_TOPIC_ARN=arn:aws:sns:us-east-1:DONDE_VA_TU_ACCOUNT_ID:critical-alerts

# OCI
OCI_COMPARTMENT_ID=ocid1.compartment.oc1..AQUI_TU_COMPARTMENT
OCI_REGION=us-chicago-1
OCI_CONFIG_FILE=/home/tu-usuario/.oci/config
OCI_CONFIG_PROFILE=DEFAULT

# Email (SES)
SES_FROM_EMAIL=alerts@sentinel-academia.com
BIENESTAR_EMAIL=bienestar@universidad-ejemplo.edu
DIRECTOR_EMAIL=director@universidad-ejemplo.edu

# Multi-tenant
DEFAULT_TENANT_ID=universidad-ejemplo

# Logging
LOG_LEVEL=INFO
POWERTOOLS_SERVICE_NAME=sentinel-academia
\end{lstlisting}

\subsection{1.2 Shared Modules}

\subsubsection{Config}

\begin{lstlisting}[language=Python, caption={src/handlers/\_shared/config.py}]
"""Configuracion tipada con pydantic-settings."""
from pydantic_settings import BaseSettings, SettingsConfigDict


class Settings(BaseSettings):
    """Settings cargados desde variables de entorno."""

    # AWS
    aws_region: str = "us-east-1"
    dynamodb_table: str = "quejas"
    s3_bucket: str = "sentinel-quejas-adjuntos"
    process_queue_url: str = ""
    critical_alerts_topic_arn: str = ""

    # OCI
    oci_compartment_id: str = ""
    oci_region: str = "us-chicago-1"

    # Email
    ses_from_email: str = "alerts@sentinel-academia.com"
    bienestar_email: str = "bienestar@universidad.edu"
    director_email: str = "director@universidad.edu"

    # Multi-tenant
    default_tenant_id: str = "universidad-ejemplo"

    # Logging
    log_level: str = "INFO"

    model_config = SettingsConfigDict(
        env_file=".env",
        case_sensitive=False,
        extra="ignore",
    )


settings = Settings()
\end{lstlisting}

\subsubsection{Logger}

\begin{lstlisting}[language=Python, caption={src/handlers/\_shared/logger.py}]
"""Logger con aws-lambda-powertools."""
from aws_lambda_powertools import Logger

logger = Logger(service="sentinel-academia")
\end{lstlisting}

\subsubsection{HTTP Response Helpers}

\begin{lstlisting}[language=Python, caption={src/handlers/\_shared/http\_response.py}]
"""Helpers para respuestas HTTP consistentes."""
import json
from typing import Any, Optional


def _response(status_code: int, body: Any, correlation_id: str) -> dict:
    return {
        "statusCode": status_code,
        "headers": {
            "Content-Type": "application/json",
            "x-correlation-id": correlation_id,
        },
        "body": json.dumps(body, default=str),
    }


def ok(body: Any, correlation_id: str) -> dict:
    return _response(200, body, correlation_id)


def accepted(body: Any, correlation_id: str) -> dict:
    return _response(202, body, correlation_id)


def bad_request(message: str, details: Optional[Any], correlation_id: str) -> dict:
    return _response(400, {
        "code": "BAD_REQUEST", "message": message, "details": details, "correlationId": correlation_id
    }, correlation_id)


def not_found(message: str, correlation_id: str) -> dict:
    return _response(404, {
        "code": "NOT_FOUND", "message": message, "correlationId": correlation_id
    }, correlation_id)


def server_error(message: str, correlation_id: str) -> dict:
    return _response(500, {
        "code": "INTERNAL_ERROR", "message": message, "correlationId": correlation_id
    }, correlation_id)
\end{lstlisting}

\subsubsection{Error Handler}

\begin{lstlisting}[language=Python, caption={src/handlers/\_shared/errors.py}]
"""Decorator para manejo uniforme de errores."""
from functools import wraps
from typing import Callable, Any
from pydantic import ValidationError
from .http_response import bad_request, server_error
from .logger import logger


def with_error_handling(handler_func: Callable) -> Callable:
    @wraps(handler_func)
    def wrapper(event: dict, context: Any) -> dict:
        try:
            return handler_func(event, context)
        except ValidationError as e:
            logger.warning("Validation error", extra={"errors": e.errors()})
            return bad_request("Validation error", e.errors(), _get_correlation_id(event))
        except Exception:
            logger.exception("Unhandled error")
            return server_error("Internal error", _get_correlation_id(event))
    return wrapper


def _get_correlation_id(event: dict) -> str:
    return event.get("headers", {}).get("x-correlation-id", "unknown")
\end{lstlisting}

\subsubsection{Correlation ID}

\begin{lstlisting}[language=Python, caption={src/handlers/\_shared/correlation\_id.py}]
"""Helper para correlation ID."""
import uuid
from typing import Optional


def get_correlation_id(event: dict) -> str:
    """Extrae correlation ID del header o genera uno nuevo."""
    headers = event.get("headers", {}) or {}
    return headers.get("x-correlation-id") or str(uuid.uuid4())
\end{lstlisting}

\subsubsection{DynamoDB Client}

\begin{lstlisting}[language=Python, caption={src/handlers/\_shared/dynamo\_client.py}]
"""Cliente DynamoDB cacheado."""
import os
from functools import lru_cache
from typing import Any
import boto3


@lru_cache(maxsize=1)
def get_dynamo_resource() -> Any:
    return boto3.resource(
        "dynamodb",
        region_name=os.environ.get("AWS_REGION", "us-east-1"),
    )


@lru_cache(maxsize=1)
def get_dynamo_client() -> Any:
    return boto3.client(
        "dynamodb",
        region_name=os.environ.get("AWS_REGION", "us-east-1"),
    )


TABLE_NAME = os.environ.get("DYNAMODB_TABLE", "quejas")
DEFAULT_TENANT_ID = os.environ.get("DEFAULT_TENANT_ID", "universidad-ejemplo")
\end{lstlisting}

\subsubsection{Pydantic Schemas}

\begin{lstlisting}[language=Python, caption={src/handlers/\_shared/schemas.py}]
"""Schemas Pydantic para input/output."""
from enum import Enum
from typing import Optional, List
from pydantic import BaseModel, Field, field_validator, ConfigDict


class CategoriaEnum(str, Enum):
    ACADEMICA = "ACADEMICA"
    INFRAESTRUCTURA = "INFRAESTRUCTURA"
    ACOSO = "ACOSO"
    ADMINISTRATIVA = "ADMINISTRATIVA"
    SALUD = "SALUD"
    OTRA = "OTRA"


class CriticidadEnum(str, Enum):
    BAJA = "BAJA"
    MEDIA = "MEDIA"
    ALTA = "ALTA"
    CRITICA = "CRITICA"


class QuejaStatusEnum(str, Enum):
    PENDIENTE = "PENDIENTE"
    EN_COLA = "EN_COLA"
    PROCESANDO = "PROCESANDO"
    ANALIZADA = "ANALIZADA"
    NOTIFICADA = "NOTIFICADA"
    ERROR = "ERROR"


class CreateQuejaInput(BaseModel):
    titulo: str = Field(min_length=5, max_length=120)
    descripcion: str = Field(min_length=20, max_length=5000)
    categoria_declarada: CategoriaEnum
    adjuntos: Optional[List[str]] = Field(default=None, max_length=5)
    anonima: bool = False
    sede: Optional[str] = Field(default=None, max_length=100)
    facultad: Optional[str] = Field(default=None, max_length=100)
    contacto_email: Optional[str] = None

    @field_validator("descripcion")
    @classmethod
    def descripcion_not_empty(cls, v: str) -> str:
        if not v.strip():
            raise ValueError("La descripcion no puede estar vacia")
        return v.strip()

    model_config = ConfigDict(extra="forbid", str_strip_whitespace=True)


class QuejaAccepted(BaseModel):
    quejaId: str
    status: QuejaStatusEnum
    correlationId: str
    createdAt: str
    estimatedAnalysisTime: Optional[int] = None


class AnalisisOutput(BaseModel):
    categoria: CategoriaEnum
    subcategoria: str
    criticidad: CriticidadEnum
    criticidadJustificacion: str
    sentimiento: str
    entidades: List[dict] = Field(default_factory=list)
    temasClave: List[str] = Field(default_factory=list)
    accionSugerida: str
    prioridad: int = Field(ge=1, le=10)
    requiereNotificacionInmediata: bool
    modeloUsado: str
    tokensConsumidos: Optional[int] = None
    latenciaMs: Optional[int] = None
\end{lstlisting}

\subsection{1.3 Lambda createQueja (la mas importante)}

\begin{lstlisting}[language=Python, caption={src/handlers/create\_queja.py}]
"""POST /api/quejas - Crea una queja y la encola para analisis."""
import json
import uuid
from datetime import datetime, timezone
from typing import Any

import boto3
from aws_lambda_powertools import Logger
from pydantic import ValidationError

from ._shared.config import settings
from ._shared.correlation_id import get_correlation_id
from ._shared.dynamo_client import get_dynamo_resource
from ._shared.errors import with_error_handling
from ._shared.http_response import accepted, bad_request
from ._shared.schemas import CreateQuejaInput

logger = Logger(service="create-queja")


@logger.inject_lambda_context(log_event=True)
@with_error_handling
def handler(event: dict, context: Any) -> dict:
    correlation_id = get_correlation_id(event)
    logger.append_keys(correlation_id=correlation_id)
    logger.info("create_queja invoked")

    # 1. Validar input
    try:
        body = CreateQuejaInput.model_validate_json(event.get("body", "{}"))
    except ValidationError as e:
        return bad_request("Invalid input", e.errors(), correlation_id)

    # 2. Generar IDs
    queja_id = str(uuid.uuid4())
    tenant_id = settings.default_tenant_id
    logger.append_keys(queja_id=queja_id, tenant_id=tenant_id)

    # 3. Persistir en DynamoDB
    table = get_dynamo_resource().Table(settings.dynamodb_table)
    now = datetime.now(timezone.utc)
    table.put_item(Item={
        "pk": f"TENANT#{tenant_id}#QUEJA#{queja_id}",
        "sk": "META",
        "tenantId": tenant_id,
        "quejaId": queja_id,
        "titulo": body.titulo,
        "descripcion": body.descripcion,
        "categoriaDeclarada": body.categoria_declarada.value,
        "status": "EN_COLA",
        "anonima": body.anonima,
        "sede": body.sede,
        "facultad": body.facultad,
        "createdAt": now.isoformat(),
        "ttl": int(now.timestamp()) + (90 * 24 * 60 * 60),
    })

    # 4. Encolar para procesamiento async
    boto3.client("sqs").send_message(
        QueueUrl=settings.process_queue_url,
        MessageBody=json.dumps({
            "eventType": "queja.creada",
            "payload": {"quejaId": queja_id, "input": body.model_dump(mode="json")},
            "metadata": {"correlationId": correlation_id},
        }),
        MessageAttributes={
            "correlationId": {"DataType": "String", "StringValue": correlation_id},
        },
    )

    logger.info("Queja creada exitosamente")
    return accepted({
        "quejaId": queja_id,
        "status": "EN_COLA",
        "correlationId": correlation_id,
        "createdAt": now.isoformat(),
        "estimatedAnalysisTime": 30,
    }, correlation_id)
\end{lstlisting}

\subsection{1.4 SAM Template (template.yaml)}

\begin{lstlisting}[language=YAML, caption={src/infra/template.yaml}]
AWSTemplateFormatVersion: '2010-09-09'
Transform: AWS::Serverless-2016-10-31

Globals:
  Function:
    Runtime: python3.12
    MemorySize: 512
    Timeout: 30
    Tracing: Active
    Environment:
      Variables:
        POWERTOOLS_SERVICE_NAME: sentinel-academia
        LOG_LEVEL: INFO
        DYNAMODB_TABLE: !Ref QuejasTable
        DEFAULT_TENANT_ID: universidad-ejemplo
        OCI_COMPARTMENT_ID: !Ref OciCompartmentId
        AWS_NODEJS_CONNECTION_REUSE_ENABLED: '1'

Parameters:
  OciCompartmentId:
    Type: String
    Description: OCI Compartment OCID for Generative AI

Resources:

  # DynamoDB Table
  QuejasTable:
    Type: AWS::DynamoDB::Table
    Properties:
      TableName: quejas
      BillingMode: PAY_PER_REQUEST
      AttributeDefinitions:
        - { AttributeName: pk, AttributeType: S }
        - { AttributeName: sk, AttributeType: S }
        - { AttributeName: quejaId, AttributeType: S }
        - { AttributeName: tenantId, AttributeType: S }
        - { AttributeName: createdAt, AttributeType: S }
        - { AttributeName: criticidad, AttributeType: S }
      KeySchema:
        - { AttributeName: pk, KeyType: HASH }
        - { AttributeName: sk, KeyType: RANGE }
      GlobalSecondaryIndexes:
        - IndexName: GSI_ById
          KeySchema: [{ AttributeName: quejaId, KeyType: HASH }]
          Projection: { ProjectionType: ALL }
        - IndexName: GSI_ByTenant
          KeySchema:
            - { AttributeName: tenantId, KeyType: HASH }
            - { AttributeName: createdAt, KeyType: RANGE }
          Projection: { ProjectionType: ALL }
      TimeToLiveSpecification:
        AttributeName: ttl
        Enabled: true

  # SQS Queues
  ComplaintsQueue:
    Type: AWS::SQS::Queue
    Properties:
      QueueName: complaints-queue
      VisibilityTimeout: 180
      MessageRetentionPeriod: 1209600
      RedrivePolicy:
        deadLetterTargetArn: !GetAtt ComplaintsDLQ.Arn
        maxReceiveCount: 3

  ComplaintsDLQ:
    Type: AWS::SQS::Queue
    Properties:
      QueueName: complaints-dlq
      MessageRetentionPeriod: 1209600

  # SNS Topic
  CriticalAlertsTopic:
    Type: AWS::SNS::Topic
    Properties:
      TopicName: critical-alerts

  # S3 Bucket
  AdjuntosBucket:
    Type: AWS::S3::Bucket
    Properties:
      CorsConfiguration:
        CorsRules:
          - AllowedOrigins: ['*']
            AllowedMethods: [GET, PUT, POST]
            AllowedHeaders: ['*']

  # API Gateway
  HttpApi:
    Type: AWS::ApiGatewayV2::Api
    Properties:
      Name: sentinel-academia-api
      ProtocolType: HTTP
      CorsConfiguration:
        AllowOrigins: ['*']
        AllowMethods: ['*']
        AllowHeaders: ['*']

  # Lambda Functions
  CreateQuejaFunction:
    Type: AWS::Serverless::Function
    Properties:
      CodeUri: ../
      Handler: handlers/create_queja.handler
      Events:
        Api:
          Type: HttpApi
          Properties:
            ApiId: !Ref HttpApi
            Path: /quejas
            Method: POST
      Policies:
        - DynamoDBCrudPolicy: { TableName: !Ref QuejasTable }
        - SQSSendMessagePolicy: { QueueName: !GetAtt ComplaintsQueue.QueueName }

  GetQuejaFunction:
    Type: AWS::Serverless::Function
    Properties:
      CodeUri: ../
      Handler: handlers/get_queja.handler
      Events:
        Api:
          Type: HttpApi
          Properties:
            ApiId: !Ref HttpApi
            Path: /quejas/{id}
            Method: GET
      Policies:
        - DynamoDBReadPolicy: { TableName: !Ref QuejasTable }

  ListQuejasFunction:
    Type: AWS::Serverless::Function
    Properties:
      CodeUri: ../
      Handler: handlers/list_quejas.handler
      Events:
        Api:
          Type: HttpApi
            Properties:
              ApiId: !Ref HttpApi
              Path: /quejas
              Method: GET
      Policies:
        - DynamoDBReadPolicy: { TableName: !Ref QuejasTable }

  ProcessQuejaFunction:
    Type: AWS::Serverless::Function
    Properties:
      CodeUri: ../
      Handler: handlers/process_queja.handler
      MemorySize: 1024
      Timeout: 300
      Events:
        SQS:
          Type: SQS
          Properties:
            Queue: !GetAtt ComplaintsQueue.Arn
            BatchSize: 5
      Policies:
        - DynamoDBCrudPolicy: { TableName: !Ref QuejasTable }
        - S3ReadPolicy: { BucketName: !Ref AdjuntosBucket }

Outputs:
  ApiUrl:
    Value: !GetAtt HttpApi.ApiEndpoint
    Export: { Name: SentinelApiUrl }
  TableName:
    Value: !Ref QuejasTable
  QueueUrl:
    Value: !Ref ComplaintsQueue
\end{lstlisting}

\subsection{1.5 Deploy del backend}

\begin{lstlisting}[language=Bash, caption={Build y deploy con SAM}]
cd src/infra

# Build
sam build
# Output: Build Succeeded

# Deploy (primera vez)
sam deploy --guided
# Responder:
#   Stack Name: sentinel-academia-dev
#   AWS Region: us-east-1
#   Parameter OciCompartmentId: <tu OCID>
#   Confirmar cambios: y
#   Permitir IAM: y
#   Save to samconfig.toml: y

# Deploys subsecuentes
sam deploy

# Ver outputs
aws cloudformation describe-stacks \
  --stack-name sentinel-academia-dev \
  --query 'Stacks[0].Outputs'
\end{lstlisting}

% ============================================
\section{FASE 2: Backend LLM (OCI Integration)}
% ============================================

\subsection{2.1 Cliente OCI Generative AI}

\begin{lstlisting}[language=Python, caption={src/services/oci\_genai.py}]
"""Cliente para OCI Generative AI Service."""
from typing import Any
import oci
from oci.generative_ai_inference import GenerativeAiInferenceClient
from oci.generative_ai_inference.models import (
    GenerateTextDetails, ChatDetails, TextContent
)


class OCIGenAIClient:
    def __init__(self, compartment_id: str, region: str = "us-chicago-1"):
        config = oci.config.from_file("~/.oci/config", "DEFAULT")
        self.client = GenerativeAiInferenceClient(config=config, region=region)
        self.compartment_id = compartment_id

    def generate_text(
        self, prompt: str, model_id: str = "cohere.command-r",
        max_tokens: int = 1000, temperature: float = 0.3
    ) -> str:
        details = GenerateTextDetails(
            compartment_id=self.compartment_id,
            serving_mode={"modelId": model_id, "servingType": "ON_DEMAND"},
            inference_request={
                "prompt": prompt,
                "maxTokens": max_tokens,
                "temperature": temperature,
                "isStream": False,
            },
        )
        response = self.client.generate_text(generate_text_details=details)
        return response.data.inference_response.inference_text
\end{lstlisting}

\subsection{2.2 LLM Router con Circuit Breaker}

\begin{lstlisting}[language=Python, caption={src/services/llm\_router.py}]
"""Router de LLM con fallback chain."""
import opossum
from .oci_genai import OCIGenAIClient

# Breakers por modelo
breaker_primary = opossum.CircuitBreaker(
    call=lambda p: oci.generate(p, "cohere.command-r"),
    timeout=30,
    error_threshold_percentage=50,
    reset_timeout=30,
)

breaker_secondary = opossum.CircuitBreaker(
    call=lambda p: oci.generate(p, "meta.llama-3.1-70b-instruct"),
    timeout=30,
    error_threshold_percentage=50,
    reset_timeout=30,
)


def generate_with_fallback(prompt: str) -> str:
    """Intenta modelos en orden, con fallback."""
    models = [
        ("cohere.command-r", breaker_primary),
        ("meta.llama-3.1-70b-instruct", breaker_secondary),
    ]

    for model_name, breaker in models:
        try:
            return breaker.call(prompt)
        except Exception:
            continue

    return json.dumps({"error": "Todos los modelos fallaron"})
\end{lstlisting}

% ============================================
\section{FASE 3: Frontend Vue 3}
% ============================================

\subsection{3.1 Setup}

\begin{lstlisting}[language=Bash, caption={Setup del frontend}]
cd web
npm install
cp .env.example .env.local
# Editar .env.local con la URL del API Gateway

npm run dev
# Abrir http://localhost:5173
\end{lstlisting}

\subsection{3.2 Deploy a Vercel}

\begin{lstlisting}[language=Bash, caption={Deploy a Vercel con CLI}]
npm install -g vercel

# Login
vercel login

# Deploy
vercel --prod

# Opciones:
# - Framework: Vite
# - Build: npm run build
# - Output: dist
# - Root: ./
\end{lstlisting}

% ============================================
\section{FASE 4: Archivos + Email}
% =========================================%

\subsection{4.1 S3 + Presigned URLs}

\begin{lstlisting}[language=Python, caption={src/handlers/get\_upload\_url.py}]
"""POST /api/uploads/presigned - Genera presigned URL para S3."""
import uuid
from typing import Any
import boto3
from aws_lambda_powertools import Logger
from botocore.config import Config
from ._shared.config import settings
from ._shared.correlation_id import get_correlation_id
from ._shared.errors import with_error_handling
from ._shared.http_response import ok

logger = Logger(service="get-upload-url")


@logger.inject_lambda_context
@with_error_handling
def handler(event: dict, context: Any) -> dict:
    correlation_id = get_correlation_id(event)
    body = json.loads(event.get("body", "{}"))
    filename = body["filename"]
    content_type = body["contentType"]
    queja_id = body["quejaId"]

    s3 = boto3.client("s3", config=Config(signature_version="s3v4"))
    key = f"tenants/{settings.default_tenant_id}/quejas/{queja_id}/{uuid.uuid4()}-{filename}"

    url = s3.generate_presigned_url(
        "put_object",
        Params={"Bucket": settings.s3_bucket, "Key": key, "ContentType": content_type},
        ExpiresIn=300,
    )

    return ok({"uploadUrl": url, "key": key, "expiresIn": 300}, correlation_id)
\end{lstlisting}

\subsection{4.2 Email con SES}

\begin{lstlisting}[language=Python, caption={src/handlers/notify.py}]
"""Consumer SNS que envia emails via SES para casos CRITICOS."""
import json
from typing import Any
import boto3
from aws_lambda_powertools import Logger
from ._shared.config import settings
from ._shared.errors import with_error_handling

logger = Logger(service="notify")


@logger.inject_lambda_context
@with_error_handling
def handler(event: dict, context: Any) -> None:
    for record in event.get("Records", []):
        sns = record["Sns"]
        message = json.loads(sns["Message"])
        analysis = message.get("detail", {}).get("analysis", {})

        if analysis.get("criticidad") != "CRITICA":
            continue

        ses = boto3.client("ses")
        ses.send_email(
            Source=settings.ses_from_email,
            Destination={"ToAddresses": [settings.bienestar_email, settings.director_email]},
            Message={
                "Subject": {"Data": f"ALERTA: Queja CRITICA {message['detail']['quejaId']}"},
                "Body": {
                    "Text": {"Data": f"Queja critica detectada.\nAccion: {analysis.get('accionSugerida', 'N/A')}"}
                },
            },
        )
        logger.info("Email sent")
\end{lstlisting}

% ============================================
\section{FASE 5: RAG (Opcional)}
% ============================================

\subsection{5.1 Ingesta del reglamento}

\begin{lstlisting}[language=Python, caption={src/services/rag\_ingest.py}]
"""Ingesta de documentos para RAG."""
import boto3
import cohere
from .oci_genai import OCIGenAIClient


def ingest_document(s3_key: str, oci_client: OCIGenAIClient):
    """Descarga documento de S3, lo chunk-ea, genera embeddings, guarda en DynamoDB."""
    s3 = boto3.client("s3")
    obj = s3.get_object(Bucket="reglamento-bucket", Key=s3_key)
    text = obj["Body"].read().decode("utf-8")

    # Chunking (500 chars con overlap de 50)
    chunks = [text[i:i+500] for i in range(0, len(text), 450)]

    # Generar embeddings con Cohere
    co = cohere.Client("tu-api-key")
    response = co.embed(texts=chunks, model="embed-multilingual-v3.0")
    embeddings = response.embeddings

    # Guardar en DynamoDB
    # ... (similar a otros repos)
\end{lstlisting}

% ============================================
\section{FASE 6: Testing}
% ============================================

\subsection{6.1 Tests unitarios}

\begin{lstlisting}[language=Bash, caption={Ejecutar tests}]
cd src
source .venv/bin/activate

# Todos los tests
pytest

# Con coverage
pytest --cov=. --cov-report=html

# Solo un test especifico
pytest tests/handlers/test_create_queja.py

# Ver reporte HTML
open htmlcov/index.html
\end{lstlisting}

\begin{lstlisting}[language=Python, caption={tests/handlers/test\_create\_queja.py}]
import json
import pytest
from moto import mock_aws
import boto3
from src.handlers.create_queja import handler


@pytest.fixture
def event():
    return {
        "body": json.dumps({
            "titulo": "Test",
            "descripcion": "Descripcion suficientemente larga",
            "categoriaDeclarada": "ACADEMICA",
        }),
        "headers": {"x-correlation-id": "test-123"},
    }


@mock_aws
def test_create_queja_success(event, monkeypatch):
    # Setup
    client = boto3.client("dynamodb", region_name="us-east-1")
    client.create_table(
        TableName="quejas",
        KeySchema=[{"AttributeName": "pk", "KeyType": "HASH"}],
        AttributeDefinitions=[{"AttributeName": "pk", "AttributeType": "S"}],
        BillingMode="PAY_PER_REQUEST",
    )
    monkeypatch.setenv("DYNAMODB_TABLE", "quejas")

    # Mock SQS
    sqs = boto3.client("sqs", region_name="us-east-1")
    q = sqs.create_queue(QueueName="test-queue")
    monkeypatch.setenv("PROCESS_QUEUE_URL", q["QueueUrl"])

    # Act
    result = handler(event, None)

    # Assert
    assert result["statusCode"] == 202
    body = json.loads(result["body"])
    assert "quejaId" in body
\end{lstlisting}

% ============================================
\section{FASE 7: Deploy Final}
% ============================================

\subsection{7.1 Deploy completo}

\begin{lstlisting}[language=Bash, caption={Deploy final paso a paso}]
# 1. Backend
cd src/infra
sam build && sam deploy

# Guardar la URL del API
API_URL=$(aws cloudformation describe-stacks --stack-name sentinel-academia-dev --query 'Stacks[0].Outputs[?OutputKey==`ApiUrl`].OutputValue' --output text)
echo "API URL: $API_URL"

# 2. Frontend
cd ../../web
# Actualizar VITE_API_URL en .env
echo "VITE_API_URL=$API_URL" > .env.production
npm run build
vercel --prod

# 3. Verificacion
curl $API_URL/quejas
# Debe responder 200
\end{lstlisting}

\subsection{7.2 Smoke tests}

\begin{lstlisting}[language=Bash, caption={Verificacion end-to-end}]
# Crear queja de prueba
curl -X POST $API_URL/quejas \
  -H "Content-Type: application/json" \
  -H "x-correlation-id: smoke-test-1" \
  -d '{
    "titulo": "Smoke test",
    "descripcion": "Verificacion de deploy",
    "categoriaDeclarada": "OTRA"
  }'

# Ver logs
sam logs -n CreateQuejaFunction --tail

# Verificar que se guardo
aws dynamodb scan --table-name quejas --max-items 1
\end{lstlisting}

% ============================================
\section{Ap\'endice A: Comandos \'utiles}
% ============================================

\begin{lstlisting}[language=Bash, caption={Cheatsheet de comandos}]
# Ver logs en tiempo real
sam logs -n CreateQuejaFunction --tail --stack-name sentinel-academia-dev

# Invocar una Lambda manualmente
sam remote invoke -n CreateQuejaFunction --event '{"body":"{}","headers":{}}'

# Ver outputs del stack
aws cloudformation describe-stacks --stack-name sentinel-academia-dev --query 'Stacks[0].Outputs'

# Borrar todo (empezar de cero)
sam delete --stack-name sentinel-academia-dev

# Ver costos del mes
aws ce get-cost-and-usage --time-period Start=2025-01-01,End=2025-01-31 --granularity MONTHLY --metrics BlendedCost

# Test del LLM con OCI
oci generative-ai-inference generate-text \
  --compartment-id $OCI_COMPARTMENT_ID \
  --serving-mode '{"modelId":"cohere.command-r","servingType":"ON_DEMAND"}' \
  --inference-request '{"prompt":"Test","maxTokens":50}'
\end{lstlisting}

% ============================================
\section{Ap\'endice B: Troubleshooting}
% ============================================

\begin{tcolorbox}[colback=red!5!white, colframe=critical, title=\textbf{LabRole does not exist}]
\textbf{Soluci\'on}: En AWS Academy, el LabRole es creado autom\'aticamente. Si no existe, ir a IAM $\rightarrow$ Roles $\rightarrow$ Create role $\rightarrow$ AWS service $\rightarrow$ Lambda $\rightarrow$ Attach policies: \texttt{AWSLambdaBasicExecutionRole} y \texttt{AWSLambdaVPCAccessExecutionRole}.
\end{tcolorbox}

\begin{tcolorbox}[colback=red!5!white, colframe=critical, title=\textbf{OCI 404 Not Found en GenAI}]
\textbf{Soluci\'on}:
\begin{enumerate}
    \item Verificar que el Compartment OCID es correcto
    \item Verificar que la regi\'on soporta GenAI (us-chicago-1, eu-frankfurt-1)
    \item Verificar policies del compartment
\end{enumerate}
\end{tcolorbox}

\begin{tcolorbox}[colback=red!5!white, colframe=critical, title=\textbf{CORS error en frontend}]
\textbf{Soluci\'on}: Verificar que el API Gateway tiene CORS configurado:
\begin{lstlisting}[language=YAML]
CorsConfiguration:
  AllowOrigins: ['*']
  AllowMethods: ['*']
  AllowHeaders: ['*']
\end{lstlisting}
\end{tcolorbox}

\begin{tcolorbox}[colback=red!5!white, colframe=critical, title=\textbf{Lambda timeout}]
\textbf{Soluci\'on}: En \texttt{template.yaml}, subir el Timeout:
\begin{lstlisting}[language=YAML]
Globals:
  Function:
    Timeout: 120  # Subir de 30 a 120 segundos
\end{lstlisting}
\end{tcolorbox}

% ============================================
\section{Ap\'endice C: Variables de entorno}
% ============================================

\begin{longtable}{|p{5cm}|p{4cm}|p{5cm}|}
\hline
\rowcolor{primary!20}
\textbf{Variable} & \textbf{Default} & \textbf{Proposito} \\
\hline
\endhead
\texttt{POWERTOOLS\_SERVICE\_NAME} & sentinel-academia & Nombre del servicio en logs \\
\texttt{LOG\_LEVEL} & INFO & Nivel de logging \\
\texttt{AWS\_REGION} & us-east-1 & Region de AWS \\
\texttt{DYNAMODB\_TABLE} & quejas & Nombre de la tabla \\
\texttt{S3\_BUCKET} & - & Bucket de adjuntos \\
\texttt{PROCESS\_QUEUE\_URL} & - & URL de SQS \\
\texttt{CRITICAL\_ALERTS\_TOPIC\_ARN} & - & ARN del topic SNS \\
\texttt{OCI\_COMPARTMENT\_ID} & - & OCID del compartment \\
\texttt{OCI\_REGION} & us-chicago-1 & Region de OCI \\
\texttt{SES\_FROM\_EMAIL} & - & Email remitente \\
\texttt{BIENESTAR\_EMAIL} & - & Email de bienestar \\
\texttt{DIRECTOR\_EMAIL} & - & Email del director \\
\texttt{DEFAULT\_TENANT\_ID} & universidad-ejemplo & Tenant por defecto \\
\end{longtable}

% ============================================
\section{Ap\'endice D: Costos esperados}
% ============================================

\begin{tabular}{|l|r|r|}
\hline
\rowcolor{primary!20}
\textbf{Servicio} & \textbf{Hackaton (4 d\'ias)} & \textbf{Producci\'on (mes)} \\
\hline
AWS Lambda & \$0.00 & \$0.01 \\
API Gateway & \$0.00 & \$0.03 \\
DynamoDB & \$0.05 & \$3.00 \\
SQS/SNS/EventBridge & \$0.00 & \$0.05 \\
S3 & \$0.00 & \$0.30 \\
CloudWatch & \$0.10 & \$0.50 \\
\textbf{Subtotal AWS} & \textbf{\$0.15} & \textbf{\$4.00} \\
\hline
OCI Cohere R & \$0.50 & \$5.00 \\
OCI Llama 3.1 & \$0.10 & \$2.00 \\
OCI Document AI & \$0.00 & \$5.00 \\
\textbf{Subtotal OCI} & \textbf{\$0.60} & \textbf{\$12.00} \\
\hline
Vercel & \$0.00 & \$0.00 \\
\hline
\textbf{TOTAL} & \textbf{\$0.75} & \textbf{\$16.00} \\
\hline
\end{tabular}

\end{document}
